\subsection{具有紧预解式的算子}

本条中，E 表示一个无限维复希尔伯特空间。

命题 19. --- 设 u 是 E 上的自伴部分算子。以下条件等价：

\begin{enumerate}
\def\labelenumi{(\roman{enumi})}
\item
  存在 E 的一个标准正交基 \(\mathrm{B} = (e_{j})_{j \in \mathrm{J}}\)，使得 u 在基 B 中是对角的，并且 u 的特征值族的绝对值沿 J 的有限子集补集滤子趋于无穷；
\item
  对每个 \(\lambda\)（属于 \(u\) 的预解集），预解式 \(\mathrm{R}(u,\lambda)\) 都是紧的；
\item
  存在复数 \(\lambda\) 属于 \(u\) 的预解集，使得预解式 \(\mathrm{R}(u,\lambda)\) 是紧的；
\item
  u 的谱与 u 的敏感谱重合。
\end{enumerate}

假设 (i) 成立，并令 \((\lambda_{j})_{j\in\mathrm{J}}\) 为 u 的特征值族。令 \(\mu\) 为 J 上的计数测度。u 的谱是该族的 \(\mu\)-本质像（IV，第 252 页命题 22），即该族 \((\lambda_{j})\) 的值集。对每个 \(\lambda\notin\mathrm{Sp}(u)\)，预解式 \(\mathrm{R}(u,\lambda)\) 在基 B 中是对角的，其特征值族为 \(((\lambda-\lambda_{j})^{-1})_{j\in\mathrm{J}}\)（同上）。由于该族趋于 0，自同态 \(\mathrm{R}(u,\lambda)\) 是紧的（IV，第 148 页命题 2(iii)）。这证明了 (i) 蕴含 (ii)。

由于 \(u\) 的预解集非空（参见 IV，第 248 页命题 17），根据 IV 第 245 页公式 (8) 和 III 第 5 页命题 3，性质 (ii) 与 (iii) 等价。

由于 \(\mathrm{R}(u,\lambda)\) 在 \(\lambda\notin\mathrm{Sp}(u)\) 时是正规的（IV，第 247 页命题 16），根据 IV 第 247 页命题 15 和 III 第 90 页命题 5，条件 (iii) 蕴含 (iv)。

最后证明 (iv) 蕴含 (i)。令 \(\mathcal{O}\) 为由 \(u\) 的特征向量组成的 E 的标准正交子集的集合。按包含关系排序的集合 \(\mathcal{O}\) 具有有限特征（E，III，第 34 页，定义 2），因为 O 属于 \(\mathcal{O}\) 当且仅当由 O 中至多两个元素组成的集合属于 \(\mathcal{O}\)。根据 E，III，第 35 页定理 1，存在极大元 O，且 O 属于 \(\mathcal{O}\)。令 F 为由 O 生成的 E 的闭子空间。对 \(e \in O\)，存在唯一的 \(\lambda(e) \in \mathbf{R}\)，使得 \(u(e) = \lambda(e)e\)（IV，第 248 页命题 17）。

 映射 \(\lambda\) 从 O 到 \(\mathbf{R}\) 的值集与算子 \(u\) 的谱重合。事实上，一方面该集合包含于算子 \(u\) 的谱；另一方面，若存在 \(\lambda_0 \in \mathrm{Sp}(u)\) 不是 \(\lambda\) 的值，则根据假设 \(\lambda_0\) 是算子 \(u\) 的特征值。存在一个范数为 1 的向量 \(e \in \mathrm{dom}(u)\)，使得 \(u(e) = \lambda_0 e\) 且 \(e \notin \mathrm{O}\)；子集 \(\mathrm{O} \cup \{e\}\) 是标准正交的（因为算子 \(u\) 的对应于不同特征值的特征空间彼此正交，这是 IV 第 248 页命题 17 的推论中的断言 \(b\)），故它属于 \(\mathcal{O}\)，这与 O 的极大性矛盾。

假设 \(F \neq E\)。则 \(F^{\circ}\) 非零。E 的自同态 \(\mathrm{R}(u,i)\) 是正规的（IV，第 247 页命题 16）。它使 E 的子空间 \(F^{\circ}\) 稳定（I，第 135 页引理 4）。令 v 为 \(F^{\circ}\) 的自同态，由 \(\mathrm{R}(u,i)\) 限制到子空间得到。它是 \(F^{\circ}\) 的连续正规自同态（同上），且其谱包含于 \(\mathrm{R}(u,i)\) 的谱中。由于 \(F^{\circ}\) 非零，v 的谱非空（I，第 26 页推论 1）。此外，v 的谱不只含 0，因为正规自同态 v 非零（I，第 110 页例 1）。取 \(s \in \mathrm{Sp}(v) - \{0\}\)。由于 s 属于 \(\mathrm{R}(u,i)\) 的谱，存在 \(e \in O\)，使得 \(s = (i - \lambda(e))^{-1}\)，且 s 是 \(\mathrm{R}(u,i)\) 的特征值（IV，第 247 页命题 15，a）。由于 s 非零，根据假设它是 \(\mathrm{Sp}(\mathrm{R}(u,i))\) 的孤立点，因而也是 \(\mathrm{Sp}(v)\) 的孤立点。因此，s 是 v 的特征值（I，第 134 页命题 5，c）。令 \(e \in F^{\circ}\) 为 v 的范数为 1 的特征向量；它也是 u 的特征向量（IV，第 247 页命题 15，b），而集合 \(O \cup \{e\}\) 与 O 在 O 中极大这一事实矛盾。因此 F = E。

因此，O 中元素组成的族是 E 的一个由 u 的特征向量组成的标准正交基。u 的谱在 R 中是闭的，敏感谱是离散的；由于这两个集合重合，对每个 R > 0，满足 \(\lambda \in \mathrm{Sp}(u)\) 且 \(|\lambda| \leqslant R\) 的元素集合是紧的，因而是有限的。因此，u 的特征值绝对值族沿 O 的有限子集补集滤子趋于无穷。故 (iv) 蕴含 (i)。

定义 12. --- 设 u 是 E 上的自伴部分算子。如果命题 19 的等价条件成立，则称 u 具有紧预解式。

命题 20. --- 设 u 是 E 上的自伴部分算子。u 具有紧预解式，当且仅当希尔伯特空间 \(E_{u}\) 到 E 的典范嵌入 j 是紧的。

假设 u 具有紧预解式。存在 \(\lambda \notin \operatorname{Sp}(u)\)，使得 \(\mathrm{R}(u, \lambda)\) 是紧的（命题 19(iii)）。令 B 为希尔伯特空间 \(E_{u}\) 的单位球。由于 u 是从 \(E_{u}\) 到 E 的连续线性映射，E 的子集 \(\mathrm{B}' = (\lambda 1_{\mathrm{E}} - u)(\mathrm{B})\) 有界。由于 \(\mathrm{R}(u, \lambda)\) 是紧的，子集 \(\mathrm{B} = \mathrm{R}(u, \lambda)(\mathrm{B}')\) 在 E 中相对紧（III，第 2 页注 1）。这证明了 j 是紧的。

 反之，假设线性映射 \(j\) 是紧的。复数 \(i\) 属于 \(u\) 的预解集（IV，第 248 页命题 17）。有 \(u \circ \mathrm{R}(u, i) = -1_{\mathrm{E}} + i\mathrm{R}(u, i)\)。令 B 为 E 的单位球。对每个 \(x \in \mathrm{B}\)，有

\[
\| u \circ \mathrm{R} (u, i) (x) \| = \| - x + i \mathrm{R} (u, i) (x) \| \leqslant 1 + \| \mathrm{R} (u, i) \|,
\]

从而

\[
\begin{array}{r l} & {\| \mathrm{R} (u, i) x \| _ {u} ^ {2} = \| \mathrm{R} (u, i) x \| ^ {2} + \| u \circ \mathrm{R} (u, i) x \| ^ {2}} \\ & {\qquad \leqslant \| \mathrm{R} (u, i) \| ^ {2} + (1 + \| \mathrm{R} (u, i) \|) ^ {2}.} \end{array}
\]

因此，B 在 \(R(u,i)\) 下的像 C 在 \(E_{u}\) 中有界。由于按假设 j 是紧的，集合 \(C = j(C)\) 在 E 中相对紧（III，同上）。所以预解式 \(R(u,i)\) 是紧的，且 u 具有紧预解式（命题 19(iii)）。

推论。 --- 设 q 是 E 上的闭正部分形式，u 是表示 q 的正自伴算子。以下条件等价：

\begin{enumerate}
\def\labelenumi{(\roman{enumi})}
\item
  部分算子 u 具有紧预解式；
\item
  E 的正自同态 \(\sqrt{(1_{\mathrm{E}} + u)^{-1}} = (1_{\mathrm{E}} + u)^{-1 / 2}\) 是紧的；
\item
  希尔伯特空间 \(E_{q}\) 到 E 的典范嵌入 j（IV，第 292 页定义 9）是紧的。
\end{enumerate}

由于 u 是正的，实数 -1 属于 u 的预解集（IV，第 248 页命题 17），故正自同态 \(v = \sqrt{(1_{\mathrm{E}} + u)^{-1}}\) 有定义，并且根据函数演算有 \(v = (1_{\mathrm{E}} + u)^{-1/2}\)。

自同态 \(v\) 是紧的，当且仅当 \(v^2 = (1_{\mathrm{E}} + u)^{-1}\) 是紧的（III，第 91 页命题 6），也就是当且仅当算子 \(u\) 具有紧预解式（命题 19(iii)）。这证明了条件 (i) 和 (ii) 等价。

假设自同态 \(v\) 是紧的。令 B 为希尔伯特空间 \(\mathrm{E}_q\) 的单位球。由于 \((1_{\mathrm{E}} + u)^{1/2}\) 是从 \(\mathrm{E}_q\) 到 E 的连续线性映射（IV，第 294 页定理 3 的推论），E 的子集 \(\mathrm{B}' = (1_{\mathrm{E}} + u)^{1/2}(\mathrm{B})\) 有界，故子集 \(j(\mathrm{B}) = \mathrm{B} = v(\mathrm{B}')\) 在 E 中相对紧（III，第 2 页注 1）。这证明了 \(j\) 是紧的。因此 (ii) 蕴含 (iii)。

线性映射 \(\widetilde{v}\colon x\mapsto(1_{\mathrm{E}}+u)^{-1/2}(x)\) 从 E 到 \(\mathrm{E}_{q}\) 定义良好且为等距映射（IV 第 294 页定理 3 的推论）。由于 \(v=j\circ\widetilde{v}\)，条件 (iii) 蕴含 (ii)（III，第 5 页命题 3）。

例。 --- 设 \(n \in N\)。令 U 为 \(R^{n}\) 的一个开子集。赋予 U 以勒贝格测度，记为 \(\mu\)。令 \(\Delta\) 为 U 上的数量微分算子 \(-\sum_{i=1}^{n}\partial_{i}^{2}\)。部分算子 \(\Delta_{-}\) 的定义域为 \(\mathcal{D}(U)\)，由 \(\varphi \mapsto \Delta(\varphi)\) 定义；它是可闭的（IV，第 242 页命题 13）且是对称的（IV，第 243 页）。它是正的，因为对每个 \(\varphi \in \mathcal{D}(U)\)，有

\[
\begin{array}{r l} & {\langle \varphi | \Delta_ {-} (\varphi) \rangle = \int_ {\mathrm{U}} \overline {{\varphi}} \Delta (\varphi) d \mu = - \sum_ {i = 1} ^ {n} \int_ {\mathrm{U}} \overline {{\varphi}} \partial_ {i} ^ {2} \varphi d \mu} \\ & {\qquad = \sum_ {i = 1} ^ {n} \int_ {\mathrm{U}} \overline {{\partial_ {i} \varphi}} \partial_ {i} \varphi d \mu = \int_ {\mathrm{U}} \sum_ {i = 1} ^ {n} | \partial_ {i} \varphi | ^ {2} d \mu \geqslant 0.} \end{array}
\]

记 \(\Delta_{D}\) 为正对称部分算子 \(\Delta_{-}\) 的弗里德里希斯扩张（IV，第 295 页，例）；它是 U 上的拉普拉斯算子，称为 U 上的狄利克雷拉普拉斯算子。

令 q 为与 \(\Delta_{D}\) 相关的正部分形式。q 的定义域是 \(\mathcal{D}(U)\) 关于正定埃尔米特形式的希尔伯特空间完备化

该形式定义为

\[
(\varphi_ {1}, \varphi_ {2}) \mapsto \int_ {\mathrm{U}} \overline {{\varphi}} _ {1} \varphi_ {2} + \sum_ {i = 1} ^ {n} \int_ {\mathrm{U}} \overline {{\partial_ {i} \varphi_ {1}}} \partial_ {i} \varphi_ {2}
\]

对所有 \((\varphi_{1},\varphi_{2})\in\mathcal{D}(\mathrm{U})\times\mathcal{D}(\mathrm{U})\) 都成立。换言之，\(q\) 的定义域是索伯列夫空间 \(\mathrm{H}_{0}^{1}(\mathrm{U})\)（IV 第 221 页第 14 号）。

*假设 U 有界。典范嵌入 \(\mathrm{H}_0^1 (\mathrm{U})\) 到 \(\mathrm{L}^2 (\mathrm{U})\) 是紧的；因此 U 上的狄利克雷拉普拉斯算子是具有紧预解式的算子（命题 20 的推论）。由于希尔伯特空间 \(\mathrm{H}_0^1 (\mathrm{U})\) 是可数型的（IV，第 222 页命题 20），且 \(\Delta_{\mathrm{D}}\) 的像是无限维的，存在一个递增的实数序列 \((\lambda_n)_{n\geqslant 0}\)，其趋于 \(+\infty\)，以及一个标准正交基 \((f_n)_{n\in \mathbf{N}}\)，它是 \(\mathrm{L}^2 (\mathrm{U})\) 的基，其元素属于 \(\Delta_{\mathrm{D}}\) 的定义域，并满足 \(\Delta_{\mathrm{D}}(f_n) = \lambda_n f_n\) 对所有 \(n\in \mathbf{N}\) 都成立。可以证明（“外尔定律”），当 T 趋于 \(+\infty\) 时，有

\[
\sum_{\substack{n\in \mathbf{N}\\ \lambda_{n}\leqslant \mathrm{T}}}1\sim \frac{c_{n}}{(2\pi)^{n}} m\mathrm{T}^{n / 2}
\]

其中 \(c_{n} = \pi^{n/2}/\Gamma(1 + n/2)\) 是 \(R^{n}\) 中单位球的体积（INT，第 V 卷，第 101 页，第 8 节第 7 号），而 \(m > 0\) 是 U 的勒贝格测度。*

